\documentclass[10pt, twocolumn]{article}
\usepackage{graphicx}
\usepackage{amsmath}
\usepackage{amssymb}
\usepackage{amsfonts}
\usepackage{textcomp}
\usepackage{amsthm}
\usepackage{mathtools}

\theoremstyle{plain}
\newtheorem{theorem}{Theorem}[section]
\newtheorem{lemma}[theorem]{Lemma}
\newtheorem{corollary}[theorem]{Corollary}
\newtheorem{proposition}[theorem]{Proposition}

\theoremstyle{definition}
\newtheorem{definition}[theorem]{Definition}
\newtheorem{example}[theorem]{Example}

\theoremstyle{remark}
\newtheorem{remark}[theorem]{Remark}

\title{%
  \vspace{-1.5em}%
  \large\textbf{Equations \& Theorems for Mathematics}%
}
\author{A. E. Mahrouss\\amlal@nekernel.org}
\date{\today}

\begin{document}

\twocolumn[{%
  \maketitle
  \vspace{-0.5em}
  \vspace{0.8em}
  \rule{\linewidth}{0.4pt}
  \vspace{1em}
}]


\section{Equations \& Theorems}

\begin{theorem}
For any domain D:
\begin{equation}
    A(n, p) = \sum_{n}^{p} M(n, p); \quad A(n, p) \in D^{+}
\end{equation}
We have:
\begin{equation}
    A(n, p) = p! + \text{ } (n+1p)! + \text{ } (n+2p)! + \text{ } \dots \text{ } + (n+Lp)! 
\end{equation}
For $A(n, p) \neq 0$. Such that:
\begin{align*}
    C &=  \frac{ p! + \text{ } (n+1p)! + \text{ } (n+2p)! + \text{ } \dots \text{ } + (n+Lp)!}{A(n, p)}\\
    A(n, p) &< \frac{\operatorname{d}}{\operatorname{dC}} < \frac{ p! + \text{ } (n+1p)! + \text{ } (n+2p)! + \text{ } \dots \text{ } + (n+Lp)!}{A(n, p)}
\end{align*}
\end{theorem}

\nocite{*}

\bibliographystyle{acm}
\bibliography{A_Course_of_Pure_Mathematics, General_Theory_Dirichlet_Series, Riemann}

\end{document}